%% =================================================================
%% Wei, Mei, Zhao, Xiao, Sun, Zhang, Guo.
%% "Nondestructively identifying the mechanical behavior of soft
%%  tissues using surface deformation with an explicit inverse
%%  approach."
%% Appl. Math. Modelling 134 (2024) 126--147.
%%
%% Reconstruction of page 4 (printed p. 129) as study notes.
%% Prose: extracted verbatim from the PDF text layer via pdftotext
%%        (line breaks and joined-word artifacts cleaned up).
%% Math:  hand-transcribed from the rendered page image — Eqs. (8)–(14)
%%        including the layered/inclusion TDF definitions, the Boolean
%%        shear-modulus composition, and the smoothed Heaviside.
%% Figure 2: placeholder box with caption + visual description.
%% Note on Eq. (12): the paper renders this single equation in bold
%%        typeface (a typesetting idiosyncrasy unique to this line).
%%        Reproduced faithfully via \boldsymbol / \mathbf.
%% =================================================================

\documentclass[11pt]{article}

\usepackage[margin=1in]{geometry}
\usepackage{amsmath, amssymb}
\usepackage{mathrsfs}
\usepackage{bm}
\usepackage{fancyhdr}

\pagestyle{fancy}
\fancyhf{}
\lhead{\small Y.~Mei et al.}
\rhead{\small Applied Mathematical Modelling 134 (2024) 126--147}
\cfoot{\small 129 $\to$ \arabic{page}}
\renewcommand{\headrulewidth}{0.4pt}

\setcounter{page}{1}
\setcounter{equation}{7}  % next equation is (8)

\begin{document}

\noindent
introducing the following topological description function (TDF):
\begin{equation}
\begin{cases}
\phi^i(\mathbf{x}) < 0 & \text{if } \mathbf{x} \in \Omega^i \\
\phi^i(\mathbf{x}) = 0 & \text{if } \mathbf{x} \in \partial\Omega^i \\
\phi^i(\mathbf{x}) > 0 & \text{if } \mathbf{x} \in \Omega \setminus \Omega^i
\end{cases}
\end{equation}

For the layered structure, $\phi^i$ is the TDF of the $i$-th
interface and $\Omega^i$ is the region occupied from the curve $C^i$
and $C^0$. The mathematical expression of TDF for a layered structure
can be defined as:
\begin{equation}
\phi^i(\mathbf{x}) = r(t) - r^i(t)
\end{equation}
where $r^i(t)$ is the distance between the curve $C^i$ and $C^0$ at
any arbitrary point on the domain. To this end, the $m$-th layer of
the structure ($m > 1$) can be defined by $\phi^{m-1}(\mathbf{x}) > 0$
and $\phi^{m}(\mathbf{x}) \leq 0$, except for the first layer
$\phi^{1}(\mathbf{x}) < 0$ and last layer $\phi^{n}(\mathbf{x}) > 0$.
To this end, A bilayer structure can be expressed by the proposed
method as shown in Fig.~2. Note that we assume that each layer is
homogeneous and did not consider spatial variations within a layer.

For the domain with inclusion embedded, $\phi^i$ is the TDF of the
$i$-th component and $\Omega^i$ is the region occupied by the $i$-th
component. For an inclusion embedded composite, the mathematical
expression of TDF is:
\begin{equation}
\phi^i(\mathbf{x}) = \sqrt{\bigl(x - x_o^i\bigr)^{2}
                           + \bigl(y - y_o^i\bigr)^{2}}
                     - r^i\!\left(\theta^i\right)
\end{equation}
where $\theta^i$ is the angle from the point $(x, y)$ to the center
of the $i$-th component $\left(x_o^i, y_o^i\right)$ in polar
coordinate system. $r^i(\theta^i)$ is the distance from the point
with the angle of $\theta^i$ on $\partial\Omega^i$ to the center
$\left(x_o^i, y_o^i\right)$ as shown in Fig.~3.

With the defined TDF, the spatial variation of the shear modulus of
a layered composite can be defined based on the Boolean operators
[31]:
\begin{equation}
\mu(\mathbf{x}) = \mu^{0}\bigl(1 - H(\phi^{1}(\mathbf{x}))\bigr)
+ \sum_{i=1}^{n-1}\!\bigl[\mu^{i}\, H(\phi^{i}(\mathbf{x}))
                           \bigl(1 - H(\phi^{i+1}(\mathbf{x}))\bigr)\bigr]
+ \mu^{n}\, H(\phi^{n}(\mathbf{x}))
\end{equation}
where $\mu^{0}$ is the shear modulus of background. $\mu^{n}$ is the
shear modulus of the last layer. $H(z)$ is the Heaviside function.
If there is no interface curve $C^i$, $\mu(\mathbf{x}) = \mu^{0}$ is
a homogeneous domain. However, the TDF definitions of composite with
embedded inclusions are different, that is:
\begin{equation}
\boldsymbol{\mu}\!\left(\boldsymbol{x}\right) =
\boldsymbol{\mu}^{0} \prod_{i=1}^{n}
  \mathbf{H}\!\left(\boldsymbol{\phi}^{i}(\boldsymbol{x})\right)
+ \sum_{i=1}^{n}\!\bigl(\boldsymbol{\mu}^{i}
  \bigl(1 - \mathbf{H}(\boldsymbol{\phi}^{i}(\boldsymbol{x}))\bigr)\bigr)
\end{equation}

As the Heaviside function is a discontinuous function, we will
utilize the following function to approximate the Heaviside function
so that we can evaluate the spatial gradient of the Heaviside
function:
\begin{equation}
H(z) = 0.5 \times \left(1 + \tanh\!\left(\dfrac{z}{\varepsilon}\right)\right)
\end{equation}
where $\varepsilon$ is a small constant. With the decrease of
$\varepsilon$, the approximation of the Heaviside function is closer
to the corresponding analytical expression (see Fig.~4). However, a
smaller $\varepsilon$ might lead to an extremely large gradient at
$z$ approaching to zero. Thus, the small constant $\varepsilon$
should be carefully chosen.

\subsection*{\textit{Inverse problem}}

\noindent
After discussing the TDF, the inverse problem can be stated as:
Given $n_{\text{meas}}$ surface measured displacement fields, find
the optimal geometric vectors $\mathbf{r}^1, \mathbf{r}^2, \ldots,
\mathbf{r}^m$ and the shear moduli $\mu^0, \mu^1, \ldots, \mu^m$,
such that the following objective function is minimized:
\begin{equation}
\Pi = \dfrac{1}{2} \sum_{i_{\text{meas}}=1}^{n_{\text{meas}}}
\int \left|\mathbf{u}_i^{\text{com}} - \mathbf{u}_i^{\text{meas}}\right|^{2}
\, d\partial\Omega_{i_{\text{meas}}} + \text{Reg}(\mu)
\end{equation}

% ---------- Figure 2 placeholder ----------
\vspace{1em}
\begin{center}
\fbox{\parbox{0.92\textwidth}{\centering\vspace{1em}
\textbf{[Figure 2 — placeholder]}\\[0.8em]
Bilayer structure illustrating the TDF construction. The bottom
layer $\mu^{0}$ (teal) is bounded below by the reference curve
$C^{0}$; the top layer $\mu^{1}$ (pink) is bounded above by the
interface curve $C^{1}$. Six control points
$\mathbf{P}^{1}_{1}, \mathbf{P}^{1}_{2}, \ldots, \mathbf{P}^{1}_{6}$
define the B-spline interface $C^{1}$ (orange). Radial distances
$r^{1}_{1}, r^{1}_{2}, \ldots, r^{1}_{6}$ connect each control point
back to $C^{0}$, parameterised along $t$.
\vspace{1em}}}
\end{center}
\vspace{0.4em}
\begin{center}
\textbf{Fig. 2.} An example of the nonhomogeneous shear modulus
distribution of a bilayer structure.
\end{center}

\end{document}
